\chapter{الشبكات الغذائية}\label{ch:g5:food-webs}

نبّهنا \cref{ch:g3:food-chains} إلى أن سلاسله الأنيقة خيوط مفردة
سُحبت من شبكة: فالصقر يأخذ الفئران أيضًا، والسحلية تأكل الذباب
أيضًا، والعشب يغذّي الأرانب كذلك. وهذه السنة نكفّ عن سحب الخيوط
وننظر إلى الشبكة كلها --- \omterm{def:g5:food-webs:web}{الشبكة الغذائية} --- وإلى ما يحفظ توازنها.

\section{من السلاسل إلى الشبكات}

\begin{definition}[الشبكة الغذائية]\label{def:g5:food-webs:web}
\emph{الشبكة الغذائية}\index{شبكة غذائية} في \omterm{def:g2:where-animals-live:habitat}{موئل} ما هي شبكة
سلاسله الغذائية \emph{كلها} في وقت واحد: فكل \omterm{def:g5:biodiversity:species}{نوع} مرسوم مرة واحدة،
ومعه سهم (وما زال يعني \emph{\omterm{def:g3:food-chains:chain}{يأكله}}،
\cref{rem:g3:food-chains:arrow}) لكل علاقة وجبة. وأكثر \omterm{def:g1:animals-around-us:animal}{الحيوانات}
تأكل أنواعًا عدة ويأكلها \omterm{def:g5:biodiversity:species}{أنواع} عدة --- فتتقاطع السلاسل ويصير الرسم
شبكة.
\end{definition}

\begin{omfigure}
\begin{tikzpicture}[scale=0.98]
  % nodes
  \node[draw, thick, omMeth, fill=omMeth!8, rounded corners=3pt,
    font=\small] (grass) at (0,0) {عشب};
  \node[draw, thick, omMeth, fill=omMeth!8, rounded corners=3pt,
    font=\small] (clover) at (2.6,0) {برسيم};
  \node[draw, thick, omMeth, fill=omMeth!8, rounded corners=3pt,
    font=\small] (oak) at (5.2,0) {أوراق بلوط};
  \node[draw, thick, omExo, fill=omExo!8, rounded corners=3pt,
    font=\small] (hopper) at (0,1.7) {جندب};
  \node[draw, thick, omExo, fill=omExo!8, rounded corners=3pt,
    font=\small] (rabbit) at (2.6,1.7) {أرنب};
  \node[draw, thick, omExo, fill=omExo!8, rounded corners=3pt,
    font=\small] (caterpillar) at (5.2,1.7) {يرقة};
  \node[draw, thick, omProp, fill=omProp!8, rounded corners=3pt,
    font=\small] (lizard) at (0,3.4) {سحلية};
  \node[draw, thick, omProp, fill=omProp!8, rounded corners=3pt,
    font=\small] (tit) at (5.2,3.4) {قرقف أزرق};
  \node[draw, thick, omThm, fill=omThm!8, rounded corners=3pt,
    font=\small] (fox) at (2.6,4.4) {ثعلب};
  \node[draw, thick, omThm, fill=omThm!8, rounded corners=3pt,
    font=\small] (buzzard) at (0.2,5.2) {صقر};
  \node[draw, thick, omThm, fill=omThm!8, rounded corners=3pt,
    font=\small] (hawk) at (5.0,5.2) {باشق};
  % arrows
  \draw[->, thick, omDef] (grass) -- (hopper);
  \draw[->, thick, omDef] (grass) -- (rabbit);
  \draw[->, thick, omDef] (clover) -- (rabbit);
  \draw[->, thick, omDef] (oak) -- (caterpillar);
  \draw[->, thick, omDef] (hopper) -- (lizard);
  \draw[->, thick, omDef] (caterpillar) -- (tit);
  \draw[->, thick, omDef] (rabbit) -- (fox);
  \draw[->, thick, omDef] (lizard) -- (buzzard);
  \draw[->, thick, omDef] (hopper) -- (tit);
  \draw[->, thick, omDef] (tit) -- (hawk);
  \draw[->, thick, omDef] (rabbit) -- (buzzard);
\end{tikzpicture}

{\small \omterm{def:g5:food-webs:web}{شبكة غذائية} مصغّرة لحافة مرج واحد. وكل سلسلة مستقيمة من
\cref{ch:g3:food-chains} موجودة فيها --- متقاطعة مع الأخريات
جميعًا.}
\end{omfigure}

\begin{example}[قراءة الشبكة]\label{ex:g5:food-webs:reading}
قراءتان للشكل. تتبّع أسهم الأرنب: فهو يأكل العشب \emph{والبرسيم}،
\omterm{def:g3:food-chains:chain}{ويأكله} الثعلب \emph{والصقر} --- أربع علاقات، حيث كانت السلسلة تُظهر
اثنتين. وعُدّ أسهم غذاء الصقر: سحلية وأرنب --- فإن فقد أحدهما بقي
الآخر. فالشبكة تسجّل بالضبط الخيارات الاحتياطية التي أثنت عليها
\cref{prop:g5:biodiversity:web}.
\end{example}

\section{الأدوار في الشبكة}

\begin{definition}[المنتِجون والمستهلِكون]\label{def:g5:food-webs:roles}
تنقسم \omterm{def:g5:biodiversity:species}{أنواع} الشبكة إلى أدوار:
\begin{itemize}
  \item \emph{المنتِجون}\index{منتج}: وهم \omterm{def:g1:plants-around-us:plant}{النباتات}، التي تصنع
        مادتها بنفسها
        (\cref{prop:g4:what-plants-need:make}) --- وكل مسار
        أسهم يبدأ عند واحد منهم
        (\cref{prop:g3:food-chains:start})؛
  \item و\emph{المستهلِكون}\index{مستهلك}: وهم \omterm{def:g1:animals-around-us:animal}{الحيوانات}، تعيش
        على غيرها --- فالعواشب تأكل المنتِجين مباشرة، واللواحم
        تأكل مستهلِكين آخرين، والقوارت تأكل الاثنين
        (\cref{def:g2:what-animals-eat:kinds})؛
  \item و\emph{المحلِّلون}\index{محلل} يعملون في الظل --- وهم
        مُعفِّنو \cref{ex:g2:caring-for-nature:compost}، يتغذون
        على البقايا الميتة ويردّون مادتها إلى التربة. وستفرد لهم
        سنة لاحقة فصلًا خاصًّا.
\end{itemize}
\end{definition}

\begin{example}[فرز الشكل]\label{ex:g5:food-webs:sorting}
في الشبكة المصغّرة: \omterm{def:g5:food-webs:roles}{المنتِجون} --- العشب والبرسيم والبلوط؛
\omterm{def:g5:food-webs:roles}{والمستهلِكون} العاشبون --- الجندب والأرنب \omterm{ex:g2:animals-grow-up:butterfly}{واليرقة}؛ \omterm{def:g5:food-webs:roles}{والمستهلِكون}
اللواحم --- السحلية والقرقف الأزرق (في أكثر غذائه) والثعلب (وهو
\omterm{def:g2:what-animals-eat:kinds}{قارت} في الحقيقة) والصقر والباشق. ثلاثة طوابق: أخضر، وآكل \omterm{def:g1:plants-around-us:plant}{نبات}،
وآكل لحم --- وكل طابق أعلى قائم على ما تحته.
\end{example}

\begin{proposition}[شكل الأعداد]\label{prop:g5:food-webs:pyramid}
عدّ سكان الشبكة طابقًا طابقًا يعطي شكلًا ثابتًا: \emph{كثير} من
\omterm{def:g5:food-webs:roles}{المنتِجين}، و\emph{أقل} من العواشب، و\emph{قليل} من اللواحم ---
فمرج فيه ملايين \omterm{def:g1:plants-around-us:plant}{نبتات} العشب يغذّي آلاف الجنادب، وهي تغذّي حفنة من
السحالي، وهي تغذّي زوجًا واحدًا من الصقور. وكل طابق يعيش على الذي
تحته ولا يستطيع أن يفوقه عددًا. ولهذا كان كبار الصائدين نادرين
دائمًا.
\end{proposition}
\admitted

\section{التوازن وما يخلّ به}

\begin{proposition}[توازن الشبكة]\label{prop:g5:food-webs:balance}
إذا تُركت الشبكة لنفسها حفظت أعدادها في \emph{توازن}\index{توازن
(الشبكة الغذائية)} تقريبي: فإن تكاثرت الأرانب أكل صائدوها أحسن
وتكاثروا بعدها، فذاب الفائض؛ وإن تضاءلت الأرانب انصرف الصائدون
الجياع إلى فريسة أخرى (الأسهم الاحتياطية) فتعافت أعداد الأرانب.
والتصحيحات في \cref{prop:g3:food-chains:depend}، إذ تجري على أسهم
كثيرة في وقت واحد، تثبّت الكل.
\end{proposition}
\admitted

\begin{example}[الخطأ الكلاسيكي]\label{ex:g5:food-webs:mistake}
حمل المزارعون يومًا حملة على إبادة الجوارح --- ``لصوص الدجاج''.
وحيث نجحوا تكاثرت الفئران والفئران الحقلية، وقد تحررت من صائديها
الأوّلين، حتى صارت آفات أكلت من الحبّ أضعاف ما كانت تساويه دجاجة
الصقر العارضة. إنها حكاية حلزون البستاني
(\cref{exo:g3:food-chains:9}) على مقياس ريف بأكمله: أزل الطابق
الأعلى من شبكة، ينفجر الطابق الذي تحته.
\end{example}

\begin{method}[التنبؤ باضطراب]\label{met:g5:food-webs:predict}
لتتنبأ بأثر إزالة \omterm{def:g5:biodiversity:species}{نوع} (أو إضافته):
\begin{enumerate}
  \item جد \omterm{def:g5:biodiversity:species}{النوع} في الشبكة وعدّد \emph{كل} أسهمه، الداخلة
        والخارجة؛
  \item فأغذيته، إذ يقلّ أكلها، تميل إلى الزيادة؛ وآكلوه، إذ
        يفتقرون، يميلون إلى النقصان أو إلى تبديل الفريسة؛
  \item وتتبّع كل تغير خطوة أخرى على الشبكة --- فآثار الخطوة
        الثانية كثيرًا ما تفاجئ؛
  \item وراجع أي الجيران لهم أسهم احتياطية (فهم يمتصّون الصدمة)
        وأيهم اعتمد على هذا \omterm{def:g5:biodiversity:species}{النوع} وحده (فهم في خطر).
\end{enumerate}
\end{method}

\section{تمارين}

\begin{exercise}[$\star$]\label{exo:g5:food-webs:1}
ما \omterm{def:g5:food-webs:web}{الشبكة الغذائية}، ولماذا تتقاطع السلاسل؟
\end{exercise}

\begin{exercise}[$\star$]\label{exo:g5:food-webs:2}
سمِّ الأدوار الثلاثة في \cref{def:g5:food-webs:roles}، مع مثال
واحد على كل واحد من الشكل (عدا المحلِّلين).
\end{exercise}

\begin{exercise}[$\star$]\label{exo:g5:food-webs:3}
من الشكل: عدّد كل ما \omterm{def:g3:food-chains:chain}{يأكله} القرقف الأزرق، وكل ما \omterm{def:g3:food-chains:chain}{يأكله}.
\end{exercise}

\begin{exercise}[$\star$]\label{exo:g5:food-webs:4}
استخرج من الشكل سلسلتين كاملتين تشتركان في الجندب.
\end{exercise}

\begin{exercise}[$\star$]\label{exo:g5:food-webs:5}
اذكر شكل الأعداد طابقًا طابقًا. ولماذا لا يستطيع طابق أن يفوق الذي
تحته عددًا؟
\end{exercise}

\begin{exercise}[$\star$]\label{exo:g5:food-webs:6}
لماذا كانت الصقور نادرة والجنادب لا تُعدّ، في جملة واحدة؟
\end{exercise}

\begin{exercise}[$\star$]\label{exo:g5:food-webs:7}
تمرّ الأرانب بسنة جيدة فتتكاثر. فماذا يحدث بعد ذلك وفق
\cref{prop:g5:food-webs:balance} --- وما الذي يذيب الفائض؟
\end{exercise}

\begin{exercise}[$\star$]\label{exo:g5:food-webs:8}
جرّب: ارسم \omterm{def:g5:food-webs:web}{الشبكة الغذائية} \omterm{def:g2:where-animals-live:habitat}{لموئل} تعرفه --- حديقة، أو متنزه، أو
بركة --- بستة \omterm{def:g5:biodiversity:species}{أنواع} على الأقل وبكل سهم تستطيع الدفاع عنه. وميّز
\omterm{def:g5:food-webs:roles}{المنتِجين} بالأخضر.
\end{exercise}

\begin{exercise}[$\star\star$]\label{exo:g5:food-webs:9}
طبّق \cref{met:g5:food-webs:predict} على إزالة السحلية من الشكل:
آثار الخطوة الأولى، ثم خطوة أخرى.
\end{exercise}

\begin{exercise}[$\star\star$]\label{exo:g5:food-webs:10}
أعد \cref{ex:g5:food-webs:mistake} سلسلةَ تنبؤات: أُزيلت الجوارح،
فماذا، ثم ماذا؟ وأي خطوة أخفق أصحاب الحملة في التنبؤ بها؟
\end{exercise}

\begin{exercise}[$\star\star\star$]\label{exo:g5:food-webs:11}
في شبكة الشكل، قارن إزالة البلوط (وهو منتِج) بإزالة الباشق (وهو
صائد أعلى). فأي الجيران يحسّ بكل فقد أولًا، وفي أي اتجاه يسري كل
اضطراب، وأي فقد تمتصّه الشبكة هنا أسهل؟ استعمل
\cref{ex:g3:food-chains:dry,ex:g3:food-chains:hunter}.
\end{exercise}
